\begin{explanationbox}
\textbf{先看一个简单模型。} $y^2=4x$ 与 $x+y-2=0$ 联立。
先写成 $x=2-y$，只需代入一次：$y^2=4(2-y)$。
\begin{center}\begin{tikzpicture}[font=\small]
\draw[black!30,rounded corners=2pt] (0,0) rectangle (14.4,8.2);
\node[anchor=west,font=\small\bfseries] at (.5,7.72) {先选纵坐标：每个实根唯一确定一个交点};
% CURVE lateral 2 1
\begin{scope}[shift={(3.8,5.1)},scale=0.6]
\begin{scope}
\clip (-1,-6.2) rectangle (8.3,2.2);
\draw[blue,very thick,domain=-6.0:2.2,samples=181,variable=\t] plot ({\t*\t/(4)},\t);
\end{scope}
\draw[-Stealth,black!45] (-1,0)--(8.3,0) node[right] {$x$};
\draw[-Stealth,black!45] (0,-6.2)--(0,2.2) node[above] {$y$};
\draw[orange!85!black,very thick] (-0.2,2.2)--(8.2,-6.2);
\coordinate (P) at (0.5358983849,1.4641016151);\fill[black] (P) circle (2pt) node[right=10pt] {$P$};
\coordinate (Q) at (7.4641016151,-5.4641016151);\fill[black] (Q) circle (2pt) node[above right=9pt] {$Q$};
\draw[teal!70!black,dashed,thick] (P)--(0,1.4641016151);\draw[teal!70!black,dashed,thick] (Q)--(0,-5.4641016151);
\node[left=9pt] at (0,1.4641016151) {$y_1$};\node[left=9pt] at (0,-5.4641016151) {$y_2$};\end{scope}\node[anchor=west] at (10.1,6.3) {$x=2-y$};\node[anchor=west] at (10.1,5.3) {$y^2+4y-8=0$};\node[anchor=west] at (10.1,4.3) {每个实根对应一点};
\node[anchor=west] at (.7,.55) {$y^2=4x,\quad x+y-2=0$；先读纵坐标根，再恢复 $x=2-y$};
\end{tikzpicture}\end{center}

两个纵坐标根为 $y=-2\pm2\sqrt3$。每个根用 $x=2-y$ 唯一恢复一个交点。
\textbf{优先选纵坐标，是因为右开口抛物线已经把 $x$ 写成了 $y$ 的函数。}
本条仍保留横坐标固定形式，便于题目只问 $x_1+x_2$、$x_1x_2$ 时速查。
\end{explanationbox}
\clearpage
\begin{explanationbox}
\textbf{投影重合，不能替代实交点判断。}
对竖直线 $x=x_0$，平方消元只得到 $(x-x_0)^2=0$。
\begin{center}\begin{tikzpicture}[font=\small]
\draw[black!30,rounded corners=2pt] (0,0) rectangle (14.4,8.2);
\node[anchor=west,font=\small\bfseries] at (.5,7.72) {竖直线的三个位置：横坐标方程都重根，交点数却不同};
% CURVE lateral 2 1
\begin{scope}[shift={(6.5,3.7)},scale=0.9]
\begin{scope}
\clip (-2,-2.5) rectangle (4,2.5);
\draw[blue,very thick,domain=-2.5:2.5,samples=181,variable=\t] plot ({\t*\t/(4)},\t);
\end{scope}
\draw[-Stealth,black!45] (-2,0)--(4,0) node[right] {$x$};
\draw[-Stealth,black!45] (0,-2.5)--(0,2.5) node[above] {$y$};
\draw[red!75!black,thick,dashed] (-1,-2.5)--(-1,2.5);
\draw[violet,very thick] (0,-2.5)--(0,2.5);
\draw[orange!85!black,very thick] (1,-2.5)--(1,2.5);
\coordinate (T) at (0.0000000000,0.0000000000);\fill[black] (T) circle (2pt);
\coordinate (P) at (1.0000000000,2.0000000000);\fill[black] (P) circle (2pt);
\coordinate (Q) at (1.0000000000,-2.0000000000);\fill[black] (Q) circle (2pt);
\end{scope}\node[anchor=west,red!75!black] at (.8,6.8) {$x=-1$：无交点};\node[anchor=west,violet] at (5.0,6.8) {$x=0$：相切};\node[anchor=west,orange!85!black] at (9.0,6.8) {$x=1$：两交点};
\node[anchor=west] at (.7,.55) {$y^2=4x$；应看 $y^2=4x_0$，不能只看 $(x-x_0)^2=0$};
\end{tikzpicture}\end{center}

$x_0=-1$ 时没有实点，$x_0=0$ 时只有顶点切点，$x_0=1$ 时有 $(1,2)$、$(1,-2)$ 两点。
三个横坐标方程的判别式都是 $0$；应回到 $y^2=4x_0$ 判断。
对于 $y^2=2spx$，合法的纵坐标二次式可用 $H$ 一次完成分类，无须把竖直弦当成例外硬记。
\end{explanationbox}
\clearpage
\begin{explanationbox}
\textbf{一个交点，也不一定是切点。}
平行于抛物线轴的直线，只有一个交点，因为代入后是一次关系。
\begin{center}\begin{tikzpicture}[font=\small]
\draw[black!30,rounded corners=2pt] (0,0) rectangle (14.4,8.2);
\node[anchor=west,font=\small\bfseries] at (.5,7.72) {一个交点有两种原因：平行于轴，或者相切};
% CURVE lateral 2 1
\begin{scope}[shift={(2.4,3.6)},scale=0.7]
\begin{scope}
\clip (-1,-2.5) rectangle (4,2.5);
\draw[blue,very thick,domain=-2.5:2.5,samples=181,variable=\t] plot ({\t*\t/(4)},\t);
\end{scope}
\draw[-Stealth,black!45] (-1,0)--(4,0) node[right] {$x$};
\draw[-Stealth,black!45] (0,-2.5)--(0,2.5) node[above] {$y$};
\draw[orange!85!black,very thick] (-1,1)--(4,1);
\coordinate (P) at (0.2500000000,1.0000000000);\fill[black] (P) circle (2pt) node[above right=9pt,xshift=12pt] {$P$};
\end{scope}% CURVE lateral 2 1
\begin{scope}[shift={(9.4,3.6)},scale=0.7]
\begin{scope}
\clip (-1,-2.5) rectangle (4,2.5);
\draw[blue,very thick,domain=-2.5:2.5,samples=181,variable=\t] plot ({\t*\t/(4)},\t);
\end{scope}
\draw[-Stealth,black!45] (-1,0)--(4,0) node[right] {$x$};
\draw[-Stealth,black!45] (0,-2.5)--(0,2.5) node[above] {$y$};
\draw[violet,very thick] (0,-2.5)--(0,2.5);
\coordinate (T) at (0.0000000000,0.0000000000);\fill[black] (T) circle (2pt) node[below left=9pt] {$T$};
\end{scope}\node at (3.6,6.7) {$y=1$：一次交点};\node at (10.6,6.7) {$x=0$：切点};
\node[anchor=west] at (.7,.55) {同一抛物线 $y^2=4x$：左图横线穿过曲线；右图在顶点相切};
\end{tikzpicture}\end{center}

左图 $y=1$ 唯一确定 $x=1/4$；右图 $x=0$ 得到 $y^2=0$，才对应二次重根。
所以先检查二次项是否消失，再讨论判别式。
上下开口时交换坐标，平行于轴的直线变成竖直线，处理逻辑相同。
\end{explanationbox}
\clearpage
